%% kapa: exam=1 round=27 year=2016 period=1 form=A subject=real_estate subject_ko=부동산학원론 seq=40 number=120
%% source: Q-Net 감정평가사 기출문제 (cst00309) artlSeq=5169402 file=extracted/1008/5169402/제27회 감정평가사 1차 문제/제27회 감정평가사 1교시문제(A형).pdf
%% conversion: kordoc 4.24.51 -> tools/gen1.py (auto); review=proofread-sonnet; flags=tikz
\begin{question}{120}
C도시 인근에 A할인점과 B할인점이 있다. 허프(D.L.Huff)의 상권분석모형을 적용할 경우, A할인점의 이용객수는 C도시 인구의 몇 \%인가? (단, 거리에 대한 소비자의 거리마찰계수값은 2이고, C도시 인구 중 50 \%가 A할인점이나 B할인점을 이용함)
\begin{center}
\begin{tikzpicture}[font=\sffamily\bfseries\scriptsize]
  \draw[thick] (0,0) circle (1.87);
  \draw[thick] (-3.5,1.25) circle (1.03);
  \draw[thick] (2.7,-0.08) circle (0.5);
  \draw (-2.62,0.72) -- (-1.78,0.55);
  \draw (1.87,-0.02) -- (2.2,-0.04);
  \node at (-2.3,0.62) [above] {8km};
  \node at (2.03,-0.02) [above] {2km};
  \node[align=center] at (0,0.35) {C도시};
  \node[align=left] at (0,-0.25) {$\cdot$ 소비자 거주지\\[2pt] $\cdot$ 인구 20만명};
  \node[align=left] at (-3.8,-0.55) {A할인점\\ $\cdot$ 면적 : 8,000m$^2$};
  \node[align=left] at (3.0,-1.0) {B할인점\\ $\cdot$ 면적 : 2,000m$^2$};
\end{tikzpicture}
\end{center}
\begin{choices}
  \item 5.0 \%
  \item 10.0 \%
  \item 15.0 \%
  \item 20.0 \%
  \item 25.0 \%
\end{choices}
\end{question}
\answer{2}
